\chapter{The atlas of languages}
\label{ch:targets}

A target language joins the atlas through two files. The table \texttt{src/lng/<target>\_lng.tsv} is its chart.
It gives each symbol of the language, its bridge word, a probe and a witness. The Lean file
\texttt{theory/lng/Lng/<target>\_lng.lean} gives a theorem for each witnessed relation, with the name of its bridge
word. C and Python are the first two targets. Each chart holds 38 rows: 29 core ops, one plain word and 8 types.

\section{The witness}

\begin{definition}
A row of $\tau_L$ with the anchor $a = \alpha(\tau_L(\sigma))$ is witnessed where the tool of $L$ runs the row's
probe on every case of $a$ and prints the answer of each case. For C the tool is a C11 compiler. For Python it is
the interpreter. A row with no witness is a hole, and so is a bridge word with no anchor.
\end{definition}

A witness is evidence on finitely many words, and it is not a proof. The Lean file of the language supplies the
proof for every word.

\section{The charts in Lean}

\texttt{harness.py lng prove} writes every row of every chart into \texttt{Lng/charts.lean}. Lean then reads the
charts as data, and \texttt{Lng/atlas.lean} proves facts about them.

\begin{theorem}[\texttt{charted}]
Each row of each chart is a type, or a witnessed relation whose bridge word has an anchor.
\end{theorem}

\begin{theorem}[\texttt{lifts}]
Each relation of each chart has a row with the same bridge word in every other chart.
\end{theorem}

The first theorem closes the gap between a chart and the anchors. The second shows that no relation of C or of
Python lacks its partner in the other chart. The kernel decides both on the 76 rows of the two charts.

\section{Fibers and cardinality}

\begin{definition}
The fiber of an anchor $a$ in the chart of $L$ is the set of symbols that reach $a$:
\[
  F_L(a) = \{\, \sigma \in \Sigma_L \;:\; \alpha(\tau_L(\sigma)) = a \,\}.
\]
A projection over $a$ from $L$ to $M$ pairs a symbol of $F_L(a)$ with a symbol of $F_M(a)$.
\end{definition}

The four cardinalities of the matrix projection read off the fibers. A projection is 1:1 where both fibers hold
one symbol. It is many:many where both hold more than one. It is 1:many where one symbol projects into the fibers
of many languages, and n:many where the fibers of $n$ languages project into those of many.

\begin{theorem}[\texttt{fibers}]
For each anchor that a chart reaches, the fiber in each chart has as many symbols as the fiber in every other chart.
\end{theorem}

\begin{theorem}[\texttt{many\_to\_many}]
The fiber of add in C holds $+$ and $+{=}$. The fiber of add in Python holds the masked $+$ and the masked $+{=}$.
\end{theorem}

\section{The semantics of a target}

For C, the operations on \texttt{uint32\_t} are the operations of $W_{32}$, since C11 6.2.5 paragraph 9 makes
unsigned arithmetic modular. Most theorems of \texttt{Lng/c\_lng.lean} hold by definition. A shift takes a count
below 32, and a quotient and a remainder take a divisor that is not 0, since C11 6.5.7 and 6.5.5 give no value
outside those ranges.

For Python, an integer has no width \cite{src:python}. The table bridges each operator through a mask, as
\texttt{(a + b) \& 0xffffffff}. Lean models the mask as the low 32 places of the two's complement of the integer.
The quotient of $/\!/$ rounds toward $-\infty$, and Lean models it with \texttt{Int.fdiv}.

\begin{theorem}[\texttt{add} of \texttt{Lng.python}]
For all $a, b \in W_{32}$, the mask of $a + b$ over the integers is $R_{\mathrm{add}}(a, b)$.
\end{theorem}

\begin{proposition}[\texttt{add\_leaves}, \texttt{shl\_count\_leaves}]
Without the mask, $a + b$ leaves $W_{32}$ for some $a$ and $b$. With the mask, $a \ll b$ differs from UP for some
count of 32 or more.
\end{proposition}

The two propositions show why each Python symbol carries a mask, and where the Python shift and the ladder part.

\section{The check}

\texttt{lng\_check} reads every target table. It reports a type or a relation that contradicts its bridge word,
and every hole. It reads the theorem names of the Lean file of each target, and it reports each witnessed relation
whose bridge word names no theorem there. A target with no Lean file is a finding too.

\section{The lift}

\begin{definition}
A chart at $n$ places and the sign reading $s$ is a set of symbols $\Sigma$, a map $\kappa : \Sigma \to A$ to the
anchors, a domain $D_\sigma \subseteq W_n$ for each symbol, and a semantics $\lBrack\cdot\rBrack$. It is sound where
$\lBrack\sigma\rBrack(x, y) = R^{s}_{\kappa(\sigma)}(x, y)$ for each $\sigma \in \Sigma$, each $x \in W_n$ and
each $y \in D_\sigma$.
\end{definition}

\texttt{Lng/any.lean} states the definition as the Lean structure \texttt{Chart n s}. Its field \texttt{sound} is
the proof, and Lean accepts no chart without it.

\begin{theorem}[\texttt{lift}]
Let $L$ and $M$ be sound charts at $n$ places and the sign reading $s$, for any $n$ and $s$. Let
$\sigma \in \Sigma_L$ and $\sigma' \in \Sigma_M$ with $\kappa_L(\sigma) = \kappa_M(\sigma')$. Then
$\lBrack\sigma\rBrack_L(x, y) = \lBrack\sigma'\rBrack_M(x, y)$ for each $x \in W_n$ and each
$y \in D_\sigma \cap D_{\sigma'}$.
\end{theorem}

\begin{proof}
Both sides equal $R^{s}_{\kappa(\sigma)}(x, y)$, and equality is transitive.
\end{proof}

The theorem asks for one anchor, and not for one bridge word. Two symbols with different bridge words lift to
each other where their bridge words reach one anchor. This is the many:many projection of
Chapter~\ref{ch:introduction}: add and add\_assign are two bridge words of one anchor.

\begin{theorem}[\texttt{lift\_all}]
Let $P$ be a finite list of symbols, each of a sound chart, that reach one anchor. Then each pair of symbols of $P$
agree on each pair of words of both domains.
\end{theorem}

The list can hold any number of charts and any number of symbols of each chart. The theorem therefore covers the
1:1, 1:many, many:many and n:many projections together.

\section{The charts proved sound}

\begin{theorem}[\texttt{g}]
For each $n$ and each sign reading $s$, the operators of \texttt{.g} form a sound chart at $n$ places, with each
domain equal to $W_n$.
\end{theorem}

\begin{theorem}[\texttt{c}, \texttt{python}]
C and Python form sound charts at 32 places, read unsigned, over their 30 shared bridge words. The domain of a
shift is a count below 32. Every other domain is $W_{32}$.
\end{theorem}

The semantics of C here is the semantics of \texttt{uint32\_t} under C11, and the semantics of Python is the
masked integer of its table. A quotient and a remainder by 0 have no value in either language or in the relation,
and the two sides agree there too.

\begin{theorem}[\texttt{many\_to\_many}]
For all $x, y \in W_{32}$, the \texttt{.g} operator $+$, the C symbols $+$ and $+{=}$, and the masked Python symbols
$+$ and $+{=}$ give one value.
\end{theorem}

\texttt{Lng/lift.lean} also states the lift from C to Python for each of the 30 shared bridge words, as an
equation on the two semantics. \texttt{Lng.lift.add} states $a + b = \mathrm{mask}(a + b)$ for all
$a, b \in W_{32}$.

\texttt{lng\_lift} computes $\sigma'$ from $\sigma$. It refuses a bridge word that the transpiler lacks, a bridge
word with no anchor, and a side with no witnessed row. The theorem \texttt{lift} gives the meaning of each lift
that it does not refuse.

\section{A new language}

A new language joins with the same two files and no change to the others. Its chart maps its symbols to bridge
words, and its Lean file proves each symbol against the relation of its anchor. That proof is the field
\texttt{sound} of one value of \texttt{Chart n s}. Every lift between the new language and each other language of
the atlas then follows from \texttt{lift}, at every width and sign that the chart states. The number of proofs grows
with the number of languages, and the number of pairs does not enter.
